\section{Discussion}

The purpose of this construction is not to replace useful empirical universal relations with more machinery. It is to separate three mechanisms that a fitted relation conflates: first-order EOS insensitivity, global integrability, and nonlinear residual curvature. If the I--Love relation is recovered from these ingredients without being imposed, the same construction can be used to ask whether additional observables contain new low-response integrable combinations. If it is not recovered, the failure remains informative because dependence on the EOS metric, observable coordinates, and stellar location is explicit.

The construction also gives a direct route to understanding apparent prior dependence. A catalog-based scatter implicitly weights EOS variations according to the population of models included in that catalog. Here the weighting appears explicitly through \(C_{\rm EOS}\). Comparing several physically reasonable covariance operators therefore tests whether a relation is robust to how function space is explored, rather than only to which named EOSs happen to have been tabulated.

The proposed density-resolved kernels provide a second diagnostic unavailable to a fitted relation alone. If a phase transition or localized stiffening generates a large transverse displacement, the responsible enthalpy or density range can be identified in the response kernel. This permits a direct comparison with earlier physical pictures based on outer-core dominance, self-similarity, and proximity to incompressibility.

Only after the static construction is validated will we enlarge the observable set to quantities such as \(\bar Q\), higher multipolar Love numbers, mode frequencies, or frequency-dependent tidal response. Any apparent new relation will trigger a fresh literature audit before a novelty claim is made.

\section{Conclusions}

We have formulated neutron-star quasi-universality as a geometric property of the relativistic EOS-to-observable map. The EOS Jacobian is combined with a declared metric on EOS function space, displacement along the stellar sequence is quotiented out, and universal relations are associated locally with low-response normal covectors. Global relations require approximate integrability, while finite residual scatter begins at second order when first-order stationarity is effective.

\todo{Replace with measured conclusions only after Stages 3--6.}
